% from ega3-6-fr.tex
\oldpage[III]{162}
Modules $\mathscr H^{-n}(f_1,\mathscr P_\bullet^{(1)})$ ou tous les Modules
$\mathscr H^{-n}(f_2,\mathscr P_\bullet^{(2)})$ soient $S$-plats. On a alors un
isomorphisme canonique de bi-$\partial$-foncteurs ($n$ parcourant $\mathbf Z$)
\[
  \mathscr H^n\bigl(f_1\times_S f_2,
    \mathscr P_\bullet^{(1)}\otimes_S\mathscr P_\bullet^{(2)}\bigr)
  \xrightarrow{\sim}
  \bigoplus_{n_1+n_2=n}
  \mathscr H^{n_1}(f_1,\mathscr P_\bullet^{(1)})
  \otimes_S
  \mathscr H^{n_2}(f_2,\mathscr P_\bullet^{(2)}).
  \tag{6.7.8.1}\label{III.6.7.8.1.fr}
\]
(« formule de Künneth »).
\end{theorem}

\begin{proof}
Cela résulte de (6.7.6) et (6.7.7).

Lorsque $S$, $Y^{(1)}$ et $Y^{(2)}$ sont affines, l'isomorphisme réciproque de
(6.7.8.1) se déduit (avec les notations de (6.7.6)) de l'homomorphisme de
bicomplexes
\[
  C^\bullet(\mathfrak U^{(1)},\mathscr P_\bullet^{(1)})
  \otimes_A
  C^\bullet(\mathfrak U^{(2)},\mathscr P_\bullet^{(2)})
  \longrightarrow
  C^\bullet(\mathfrak U,\mathscr Q_\bullet)
\]
par le procédé défini dans (G, I, 2.7), comme il résulte de (G, I, 5.5).
\end{proof}

\begin{proposition}[6.7.9]
Supposons vérifiées les trois conditions suivantes : $1^\circ$ $S$,
$Y^{(1)}$ et $Y^{(2)}$ sont localement noethériens, $f_1$ et $f_2$ sont
propres, $Y^{(1)}$ ou $Y^{(2)}$ de type fini sur $S$; $2^\circ$
$\mathscr P_\bullet^{(1)}$ et $\mathscr P_\bullet^{(2)}$ sont limités
inférieurement; $3^\circ$ pour tout $n\in\mathbf Z$,
$\mathscr H_n(\mathscr P_\bullet^{(i)})$ est un Module cohérent ($i=1,2$).
Dans ces conditions,
$\mathfrak{Tor}_n^S
(f_1,f_2;\mathscr P_\bullet^{(1)},\mathscr P_\bullet^{(2)})$ est un
$\mathscr O_Y$-Module cohérent (avec
$Y=Y^{(1)}\times_S Y^{(2)}$).
\end{proposition}

\begin{proof}
Il résulte de (6.5.13) que les hypertor locaux
$\mathscr{T}\!or_n^S(\mathscr P_\bullet^{(1)},\mathscr P_\bullet^{(2)})$ sont
des $\mathscr O_X$-Modules cohérents
($X=X^{(1)}\times_S X^{(2)}$ étant localement noethérien, car un des
$X^{(i)}$ est par hypothèse de type fini sur $S$ (I, 6.3.4 et 6.3.8)). Comme
$Y$ est localement noethérien et $f_1\times_S f_2$ propre (II, 5.4.2), il
résulte de (6.2.5) que les termes ${}^{(b)}\mathscr E_{pq}^2$ de (6.7.3) sont
des $\mathscr O_Y$-Modules cohérents. Comme toutes les suites spectrales de
(6.7.3) sont birégulières en vertu de l'hypothèse $2^\circ$, on conclut par
($0_{\mathrm I}$, 11.1.8).
\end{proof}

\begin{env}[6.7.10]
Soient maintenant $Y'^{(i)}$ deux $S$-préschémas,
$v_i:Y'^{(i)}\to Y^{(i)}$ deux $S$-morphismes ($i=1,2$),
$v=v_1\times_S v_2$ leur produit, qui est un $S$-morphisme $Y'\to Y$, où
l'on pose $Y'=Y'^{(1)}\times_S Y'^{(2)}$. Considérons d'autre part, pour
$i=1,2$, un $S$-préschéma $X'^{(i)}$, et deux $S$-morphismes
$u_i:X'^{(i)}\to X^{(i)}$, $f'_i:X'^{(i)}\to Y'^{(i)}$ de sorte que les
diagrammes
\[
\xymatrix@C=36pt@R=28pt{
  X'^{(i)}\ar[r]^{u_i}\ar[d]_{f'_i}
  & X^{(i)}\ar[d]^{f_i}
  \\
  Y'^{(i)}\ar[r]_{v_i}
  & Y^{(i)}
}
\tag{6.7.10.1}\label{III.6.7.10.1.fr}
\]
soient commutatifs, les morphismes $f'_i$ étant \emph{séparés et
quasi-compacts}. On a alors des $\mathscr O_{Y'}$-homomorphismes canoniques de
foncteurs spectraux
\[
  v^*\bigl({}^{(t)}\mathscr E
  (f_1,f_2;\mathscr P_\bullet^{(1)},\mathscr P_\bullet^{(2)})\bigr)
  \longrightarrow
  {}^{(t)}\mathscr E
  (f'_1,f'_2;u_1^*(\mathscr P_\bullet^{(1)}),
                 u_2^*(\mathscr P_\bullet^{(2)}))
  \tag{6.7.10.2}\label{III.6.7.10.2.fr}
\]
pour $t=a$, $a'$, $b$, $b'$, $c$, $d$. Pour les définir, supposons d'abord
$S=\operatorname{Spec}(A)$,
$Y^{(i)}=\operatorname{Spec}(B_i)$,
$Y'^{(i)}=\operatorname{Spec}(B'_i)$ affines; les $X^{(i)}$ et $X'^{(i)}$
sont alors des schémas quasi-compacts. Pour calculer les suites spectrales
${}^{(t)}\mathscr E
(f_1,f_2;\mathscr P_\bullet^{(1)},\mathscr P_\bullet^{(2)})$, nous
considérons comme dans (6.6.1) des recouvrements finis
$\mathfrak U^{(i)}$ par des ouverts affines de $X^{(i)}$ ($i=1,2$); pour
calculer
${}^{(t)}\mathscr E
(f'_1,f'_2;u_1^*(\mathscr P_\bullet^{(1)}),
u_2^*(\mathscr P_\bullet^{(2)}))$, nous considérons des recouvrements finis
$\mathfrak U'^{(i)}$

